% Section I. Target: ~25% of the body (see SPEC.md).
% Evidence: research/comparison.csv, research/limitations-evidence.md,
% research/peru-energy-stats.md (abstract-level; verified 2026-10-28).
\section{Introducción}
\label{sec:introduction}

% A. Problem context.
La gestión eficiente de la energía eléctrica requiere información oportuna y detallada sobre cómo, cuándo y en qué cantidad se consume.
En el Perú, la producción total de electricidad alcanzó \SI{63565}{GWh} en 2024, con un incremento de \SI{2,4}{\percent} respecto al año anterior~\cite{minem2025produccion}.
Asimismo, en el primer trimestre de 2024 el \SI{96,0}{\percent} de la población tenía acceso a la energía eléctrica por red pública~\cite{inei2024acceso}.
Con una cobertura tan amplia, el reto se desplaza progresivamente desde el acceso hacia el uso informado y eficiente de la energía.

Sin embargo, la supervisión del consumo suele basarse en lecturas periódicas o en la facturación mensual acumulada.
Estos registros agregados no permiten identificar oportunamente variaciones de demanda, picos de consumo ni patrones horarios de utilización.
Los sistemas basados en Internet de las Cosas (IoT) permiten medir variables eléctricas con mayor frecuencia, transmitirlas a plataformas digitales y analizarlas de forma continua~\cite{rind2023perspective}.
Estas capacidades resultan útiles en viviendas, edificios, laboratorios y pequeñas instalaciones industriales.
No obstante, la utilidad de un sistema de este tipo depende de la exactitud de sus mediciones, de la calidad de los datos, de la latencia de transmisión y de su capacidad para producir indicadores energéticos interpretables.

% B. Related work.
La literatura reciente presenta numerosos medidores IoT de bajo costo.
Sousa \textit{et~al.}~\cite{sousa2023smartmeter} desarrollaron un medidor inalámbrico basado en ESP32 con medición de calidad de energía para aplicaciones de redes inteligentes.
Garcés \textit{et~al.}~\cite{garces2025iotmeter} propusieron un medidor con sensores no invasivos y calibración manual de señales para entornos residenciales e industriales.
Munoz \textit{et~al.}~\cite{munoz2024smartmeter} construyeron un medidor trifásico con un costo inferior a \SI{80}{USD} que mide, entre otras magnitudes, potencias activa, reactiva y aparente, factor de potencia y armónicos, con un error absoluto menor al \SI{1}{\percent}.
Otros trabajos emplean módulos de medición integrados: Rohman \textit{et~al.}~\cite{rohman2025esp32} supervisaron generadores trifásicos con un módulo PZEM-016 y un servidor web local en el ESP32 que no requiere acceso a internet.

Un segundo grupo de trabajos enfatiza la validación experimental frente a instrumentos de referencia.
Kizonde \textit{et~al.}~\cite{kizonde2023monitoringnode} compararon un nodo de monitoreo con una unidad de medición de potencia y reportaron baja linealidad para cargas de baja potencia y un sesgo considerable para cargas de alta potencia.
Cárdenas Magaña \textit{et~al.}~\cite{cardenas2026residential} validaron un sistema ESP32 con PZEM-004T frente a un instrumento comercial calibrado, usando cargas inductivas, resistivas y electrónicas, con errores absolutos medios de \SI{0,31}{V}, \SI{0,008}{A} y \SI{1,16}{W}.
Marin \textit{et~al.}~\cite{marin2026rmsbenchmark} contrastaron mediciones RMS de tensión y corriente obtenidas con nodos Arduino y sensores de efecto Hall frente a un analizador de calidad de energía de referencia.

Un tercer grupo aborda la comunicación y la plataforma de datos.
Khan y Pirak~\cite{khan2021mqtt} evaluaron experimentalmente los protocolos MQTT y CoAP para un medidor NB-IoT, incluyendo la pérdida de paquetes en distintos entornos de instalación.
Nizam \textit{et~al.}~\cite{nizam2024realtime} integraron Modbus, ESP-NOW y MQTT en un sistema de monitoreo con actualizaciones cada \SI{5}{s}.
De Sousa \textit{et~al.}~\cite{desousa2026smartmeter} implementaron un medidor monofásico con ESP32-S3 y una plataforma basada en MQTT, InfluxDB y Grafana; reportaron un error porcentual absoluto medio de \SI{4,56}{\percent} en corriente frente a un analizador de referencia y latencias de \SIrange{203,63}{290,86}{ms} en el servidor.
Finalmente, en el análisis de datos de consumo, Oprea \textit{et~al.}~\cite{oprea2021anomaly} detectaron anomalías en datos de medidores inteligentes con una exactitud del \SI{90}{\percent}, y Zheng \textit{et~al.}~\cite{zheng2023thresholds} propusieron umbrales dinámicos para detectar fluctuaciones y picos anómalos en equipos industriales.

% C. Limitations of prior work (narrowed to what the evidence supports).
Estos antecedentes muestran avances significativos, pero también limitaciones recurrentes.
Primero, la validación metrológica es heterogénea: algunos trabajos no identifican el instrumento de referencia~\cite{rohman2025esp32}, otros reportan el error de una sola magnitud~\cite{desousa2026smartmeter} y otros evalúan un conjunto reducido de cargas~\cite{rohman2025esp32}.
Segundo, la latencia extremo a extremo rara vez se cuantifica; entre los trabajos revisados, solo~\cite{desousa2026smartmeter} reporta valores medidos.
Tercero, aunque algunos sistemas almacenan datos sin conexión~\cite{kizonde2023monitoringnode}, en los trabajos revisados no se reporta una evaluación controlada de la recuperación de mediciones tras una interrupción de la comunicación.
Cuarto, los métodos de análisis de consumo suelen evaluarse sobre conjuntos de datos independientes del hardware de medición~\cite{oprea2021anomaly,zheng2023thresholds}, por lo que su desempeño no se vincula con la calidad de las mediciones que los alimentan.

% D. Research gap.
En consecuencia, la exactitud de medición, el desempeño de la comunicación, la tolerancia a desconexiones y el análisis del consumo suelen estudiarse en trabajos separados.
En la literatura revisada no se identificó un prototipo de bajo costo que evalúe conjuntamente la exactitud frente a un instrumento de referencia, la latencia de actualización, la recuperación ante interrupciones de red y la caracterización de patrones de consumo.

% E. Objective and contributions (aligned with spec 001; no results stated).
El objetivo de este trabajo es diseñar, implementar y evaluar experimentalmente un sistema IoT para el monitoreo y análisis del consumo de energía eléctrica en tiempo real flexible, aplicado a un circuito monofásico de laboratorio, y evaluar su desempeño frente a un instrumento de referencia.
Las principales contribuciones son:
\begin{enumerate}
  \item una arquitectura IoT que integra adquisición simultánea de tensión y corriente, cálculo local de potencia activa y energía, transmisión MQTT con confirmación de entrega y cola persistente, y almacenamiento en una base de datos de series temporales;
  \item un módulo de análisis que agrega el consumo por hora, día y semana, identifica los periodos de máxima demanda y genera alertas por umbral con confirmación temporal para descartar transitorios de arranque;
  \item una evaluación experimental de la exactitud de tensión, corriente, potencia activa y energía frente a un instrumento de referencia, con cargas resistivas de distinta potencia y una carga inductiva o no lineal;
  \item una evaluación de la latencia de actualización, la confiabilidad de entrega y la recuperación de datos tras interrupciones controladas de la red; y
  \item un análisis de las limitaciones técnicas del sistema en distintas condiciones de carga.
\end{enumerate}

El resto del artículo se organiza como sigue.
La Sección~\ref{sec:methodology} describe la metodología, la arquitectura y el diseño experimental.
La Sección~\ref{sec:results} presenta y discute los resultados.
La Sección~\ref{sec:conclusions} resume las conclusiones y el trabajo futuro.
